\chapter{Resultados verificados}

\section{Línea base de gateo con cadencia corregida}

La nueva línea base de Gazebo completó 13 ciclos y su repetición independiente
otros 13, sin activaciones del supervisor. El ciclo medio de la primera
ejecución fue 4,320013~s. Al comparar los ciclos 2--13 de ambas ejecuciones, las
diferencias relativas de duración, avance, velocidad, excursión lateral,
altura, roll, pitch y salto articular permanecieron por debajo del 0,2~\%.

En régimen permanente, el primer ensayo obtuvo un avance medio de
0,023338~m/ciclo, velocidad de 0,005402~m/s, excursión lateral de 0,015133~m,
roll máximo de 2,233112 grados y pitch máximo de 4,365442 grados. Estos datos
demostraron repetibilidad numérica en simulación, no equivalencia estadística
con el robot físico.

La Figura~\ref{fig:gateo-series} presenta las series temporales de la primera
línea base corregida. El crecimiento de $x$ evidencia avance, mientras que las
oscilaciones periódicas de $y$, roll y pitch muestran la respuesta asociada a
las transferencias de peso. La Figura~\ref{fig:gateo-repro} contrasta los dos
ensayos ciclo a ciclo; el ciclo 1 se conserva y exhibe el transitorio de inicio.

\begin{figure}[p]
\centering
\includegraphics[width=0.94\textwidth]{Resultados/gateo_series_temporales.png}
\caption{Series temporales de la línea base de gateo con cadencia corregida.
Las líneas verticales delimitan ciclos completos. Fuente: elaboración propia a
partir de la bolsa \texttt{linea\_base\_cadencia\_corregida\_20260814\_1049}.}
\label{fig:gateo-series}
\end{figure}

\begin{figure}[p]
\centering
\includegraphics[width=0.96\textwidth]{Resultados/gateo_reproducibilidad.png}
\caption{Comparación ciclo a ciclo de dos ensayos independientes de gateo con
cadencia corregida. La similitud corresponde a simulación bajo configuraciones
equivalentes y no implica validación física. Fuente: elaboración propia.}
\label{fig:gateo-repro}
\end{figure}

\section{Validación de la marcha paso}

Dos ensayos independientes en Gazebo completaron 12 ciclos cada uno, con cero
eventos del supervisor y diferencias de medias inferiores al 0,4~\%. En los
ciclos 2--12, el avance medio fue 0,022031~m/ciclo, la velocidad
0,003825~m/s, el roll máximo 1,281777 grados y el pitch máximo 2,483731 grados.

MuJoCo ejecutó 12 ciclos con duración media de 5,759999~s, error RMS articular
de 0,026232~rad y error máximo de 0,054983~rad. El adaptador utilizado no
publicó una pose corporal equivalente a la de Gazebo; por tanto, esta prueba
solo respaldó cadencia y seguimiento articular simulado.

Las series de la Figura~\ref{fig:paso-series} muestran el avance y la respuesta
periódica del cuerpo durante la marcha paso. La comparación de la
Figura~\ref{fig:paso-repro} permite observar el transitorio del primer ciclo y
la proximidad de los dos ensayos posteriores. La Figura~\ref{fig:paso-mujoco}
resume el error articular y el error de duración respecto de 5,76~s en MuJoCo.

\begin{figure}[p]
\centering
\includegraphics[width=0.94\textwidth]{Resultados/paso_series_temporales.png}
\caption{Series temporales de la línea base de marcha paso en Gazebo. Fuente:
elaboración propia a partir de la bolsa \texttt{paso\_linea\_base\_20260814}.}
\label{fig:paso-series}
\end{figure}

\begin{figure}[p]
\centering
\includegraphics[width=0.96\textwidth]{Resultados/paso_reproducibilidad.png}
\caption{Reproducibilidad de la marcha paso sobre 12 ciclos equivalentes en
Gazebo. Cada curva corresponde a un ensayo independiente. Fuente: elaboración
propia.}
\label{fig:paso-repro}
\end{figure}

\begin{figure}[p]
\centering
\includegraphics[width=0.94\textwidth]{Resultados/paso_mujoco_seguimiento.png}
\caption{Seguimiento articular y error de cadencia de la marcha paso durante 12
ciclos en MuJoCo. Esta figura no contiene pose ni estabilidad corporal. Fuente:
elaboración propia a partir de \texttt{metricas\_mujoco\_por\_ciclo.csv}.}
\label{fig:paso-mujoco}
\end{figure}

\section{Contacto previsto frente a contacto observado}

La primera caracterización extensa de contacto incluyó 76 ciclos. La
coincidencia simultánea fue 32,550~\%. Las patas delanteras despegaron unos
0,38~s tarde y aterrizaron cerca de 1,364~s tarde; las patas traseras no
mostraron despegue y deslizaron. La incorporación posterior de subfases redujo
el retardo crudo de despegue, pero generó recontactos traseros breves durante el
ascenso.

Para distinguir esos eventos, el ensayo nominal del 1 de septiembre de 2026
registró 24 ciclos completos, 105,522328~s entre marcadores y 79.790 mensajes.
La Tabla~\ref{tab:contacto-crudo-filtrado} resume los episodios sin contacto.

\begin{table}[H]
\centering
\caption{Comparación entre contacto crudo y filtrado durante 24 ciclos de gateo nominal.}
\label{tab:contacto-crudo-filtrado}
\begin{tabular}{lrrrrr}
\toprule
Pata & Cruda (\%) & Filtrada (\%) & Media (s) & Máxima (s) & $\geq 0,12$ s \\
\midrule
FL & 62,304 & 60,747 & 0,944766 & 0,999304 & 49/49 \\
FR & 62,731 & 61,290 & 0,935656 & 0,994892 & 48/49 \\
RL & 89,414 & 87,717 & 0,074645 & 0,089805 & 0/24 \\
RR & 89,088 & 87,339 & 0,073803 & 0,089396 & 0/25 \\
\bottomrule
\end{tabular}
\end{table}

La coincidencia simultánea fue 20,638810~\% para el estado crudo y
13,621271~\% para el filtrado. Ninguna pérdida trasera persistió los 0,12~s
exigidos después de la declaración cruda. En consecuencia, no se obtuvieron
pares filtrados de despegue y aterrizaje para RL o RR.

La Figura~\ref{fig:contacto-persistencia} hace visible la diferencia entre las
patas delanteras y traseras. FL y FR superaron ampliamente el umbral de
persistencia, mientras que los máximos de RL y RR permanecieron debajo de la
línea de 0,12~s.

\begin{figure}[p]
\centering
\includegraphics[width=0.96\textwidth]{Resultados/contacto_persistencia_cruda.png}
\caption{Duración media y máxima de los episodios crudos sin contacto. La línea
discontinua representa la persistencia adicional exigida para confirmar una
pérdida filtrada. Fuente: elaboración propia a partir de 24 ciclos del ensayo
\texttt{contactos\_debounce\_nominal\_valido\_20260901\_0828}.}
\label{fig:contacto-persistencia}
\end{figure}

Durante esta ventana no hubo activaciones del supervisor. El ciclo medio duró
4,319990~s, el avance fue 0,023558~m/ciclo y la velocidad 0,005453~m/s. Los
máximos medios de roll y pitch fueron 2,074511 y 4,102216 grados.

\section{Calidad y reproducibilidad de los datos}

Antes de cada ensayo se auditó el grafo ROS~2 para evitar nodos duplicados. Las
bolsas sin orden inicial, sin fases o con diagnósticos saturados se marcaron
como inválidas. Los analizadores produjeron CSV por ciclo, transiciones y
episodios, además de informes regenerables. El SHA-256 completo se conservó en
el registro de ensayo; su forma abreviada fue
\texttt{80cf3361...864f82}.
